% Generated draft: requires mathematical and source review.
\documentclass[11pt]{article}
\usepackage[margin=25mm]{geometry}
\usepackage{fontspec}
\setmainfont{Latin Modern Roman}
\usepackage{amsmath,amssymb,mathtools,unicode-math}
\setmathfont{Latin Modern Math}
\usepackage{enumitem,needspace,xurl,hyperref}
\setlength{\parindent}{0pt}
\setlength{\parskip}{5pt}
\setlength{\emergencystretch}{3em}
\newcommand{\N}{\mathbb{N}}
\newcommand{\Z}{\mathbb{Z}}
\newcommand{\Q}{\mathbb{Q}}
\newcommand{\R}{\mathbb{R}}
\newcommand{\C}{\mathbb{C}}
\newcommand{\F}{\mathbb{F}}
\newcommand{\A}{\mathbb{A}}
\newcommand{\PP}{\mathbb{P}}
\newcommand{\E}{\mathbb{E}}
\newcommand{\eps}{\varepsilon}
\newcommand{\dd}{\,\mathrm{d}}
\begin{document}
\section*{UnsolvedMath ID: 701302}
\textbf{Source problem code:} catalogue entry
\subsection{Definitions and mathematical statement}
A group \(G\) is periodic if every element of \(G\) has finite order. It is simple if it is nontrivial and its only normal subgroups are \(\{1\}\) and \(G\). A \(2\)-subgroup is a subgroup all of whose elements have order a power of \(2\); here a Sylow \(2\)-subgroup means a \(2\)-subgroup maximal under inclusion.

A field \(P\) is locally finite if every subfield generated by finitely many elements is finite. Interpret the conventional notation \(L_2(P)\) as
\[
L_2(P)=\operatorname{PSL}_2(P)
=\operatorname{SL}_2(P)/Z(\operatorname{SL}_2(P)),
\]
where \(\operatorname{SL}_2(P)\) is the group of determinant-one \(2\times2\) matrices over \(P\), and \(Z\) denotes the center.

\textbf{Question (V. P. Shunkov, 2006; Kourovka Notebook 16.111).} Let \(G\) be an infinite simple periodic group possessing a dihedral Sylow \(2\)-subgroup. Must there exist a locally finite field \(P\) of odd characteristic such that
\[
G\cong\operatorname{PSL}_2(P)?
\]
Any such \(P\) must be infinite, since \(G\) is infinite.

\textbf{Unresolved qualification.} The supplied record does not define ``dihedral'' in this infinite-group setting. It does not specify whether the Sylow subgroup must be a finite dihedral \(2\)-group or whether an infinite locally dihedral \(2\)-group is also intended, nor does it specify a minimum allowed finite order. The question above retains the source terminology; its precise range requires confirmation from the source's conventions.
\subsection{Short English statement}
Shunkov asks whether every infinite simple group in which every element has finite order, and which has a dihedral Sylow \(2\)-subgroup, is isomorphic to \(\operatorname{PSL}_2(P)\) over a locally finite field of odd characteristic. The intended scope of ``dihedral'' remains to be clarified.
\subsection{Sources}
\begin{itemize}
\item Local UnsolvedMath registry, permanent ID \(701302\), exact source code \texttt{catalogue entry}, title \emph{Kourovka Notebook Problem 16.111}; supplied upstream revision \texttt{d371ee917ac9647fbb06cd6fc529deec6f3dda82}.
\item \emph{The Kourovka Notebook (Unsolved Problems in Group Theory)}, original problem 16.111, 16th issue (2006), proposed by V. P. Shunkov. The supplied background also identifies the 21st issue edition (2026), \texttt{arXiv:1401.0300v48}. The collection description names E. I. Khukhro and V. D. Mazurov as editors.
\item The supplied registry labels this problem \texttt{open}. This is a supplied status label, not an independently verified assessment of the literature. The cited source has not been independently consulted for this draft.
\end{itemize}
\end{document}
